\subsection{Remarks on Fréchet spaces}

We are going to consider prop. 2 of GT, IX, § 3.1 in the case of locally convex spaces.

PROPOSITION 11. --- Let E be a metrisable locally convex space. The topology of E can be defined by a distance that is invariant under translations, and for which the open balls are convex.

Let \((p_{n})_{n\in\mathbf{N}}\) be a sequence of semi-norms that define the topology of E. Let \(d_{n}\) be the pseudometric defined by \(d_{n}(x,y)=\inf(p_{n}(x-y),1/n)\) for x,y in E; it is invariant under translations. For every \(n\geqslant0\) , and every real number \(R\geqslant0\) , let \(B_{n,R}\) be the set of \(x\in E\) for which \(d_{n}(x,0)<R\) . If \(R\geqslant1/n\) , then \(B_{n,R}=E\) , and in the other case \(B_{n,R}\) is formed from the \(x\in E\) such that \(p_{n}(x)<R\) ; in all cases \(B_{n,R}\) is convex.

For \(x, y\) in \(\mathbf{E}\) define \(d(x, y) = \sup_{n \in \mathbb{N}} d_n(x, y)\) . We see immediately that \(d\) is a distance, invariant under translations on \(\mathbf{E}\) and defining the topology of \(\mathbf{E}\) . For \(x_0 \in \mathbf{E}\) and \(\mathbf{R} \geqslant 0\) , the open ball with centre \(x_0\) and radius \(\mathbf{R}\) (for the distance \(d\) ) is equal to \(\bigcap_{n \in \mathbb{N}} (x_0 + B_{n,\mathbf{R}})\) , therefore it is convex.

PROPOSITION 12. --- Let E and F be two Fréchet spaces and u a continuous linear mapping of E on F. Then there exists a section of u that is continuous though not necessarily linear.

By prop. 11 there exists a distance d in E, invariant under translations, defining the topology of E and for which open balls are convex. Given y and \(y'\) in F, let \(\delta(y, y')\) be the distance apart of the closed sets \(u^{-1}(y)\) and \(u^{-1}(y')\) in E. As u is a strict morphism (I, p.~17, th. 1) the remark of GT, IX, §3.1 shows that \(\delta\) is a distance on F defining the topology of F. We shall construct, inductively, a sequence of continuous mappings \((s_n)_{n \in \mathbf{N}}\) of F in E satisfying the following inequalities for all \(y \in F\) :

\[
\delta (y, u (s _ {n} (y))) <   2 ^ {- n}\tag{2}
\]

\[
d \left(s _ {n} (y), s _ {n - 1} (y)\right) <   2 ^ {- n + 1} \quad (\text { only   if } n \geqslant 1).\tag{3}
\]

Suppose then that either \(n = 0\) , or \(n \geqslant 1\) and that \(s_{n-1}\) has been constructed. Let \(y_0 \in \mathbf{F}\) ; as \(u\) is surjective, the set \(u^{-1}(y_0)\) is non-empty, and for \(n \geqslant 1\) , we have \(d(u^{-1}(y_0), s_{n-1}(y_0)) < 2^{-n+1}\) by the induction hypothesis. Therefore there exists a point \(x_0\) of \(\mathbf{E}\) such that \(u(x_0) = y_0\) and for \(n \geqslant 1\) , \(d(x_0, s_{n-1}(y_0)) < 2^{-n+1}\) . As the mapping \(s_{n-1}\) is continuous, the set of points \(y\) of \(\mathbf{F}\) which satisfy the inequalities \(\delta(y, y_0) < 2^{-n}\) and \(d(x_0, s_{n-1}(y)) < 2^{-n+1}\) is an open neighbourhood of \(y_0\) . Hence there exist an open covering \((\mathrm{V}_i)_{i \in \mathbf{I}}\) of \(\mathbf{F}\) and constant mappings \(s_{n,i}\) of \(\mathbf{F}\) in \(\mathbf{E}\) which satisfy the inequalities (2) and (3) in \(\mathrm{V}_i\) where one replaces \(s_n\) by \(s_{n,i}\) . As the space \(\mathbf{F}\) is metrisable, there exists a continuous partition of unity \((f_i)_{i \in \mathbf{I}}\) , that is locally finite and subordinate to the covering \((\mathrm{V}_i)_{i \in \mathbf{I}}\) (GT, IX, § 4.5, th. 4 and § 4.4, cor. 1). For every \(y \in \mathbf{F}\) , put \(s_n(y) = \sum_{i \in \mathbf{I}} f_i(y) \cdot s_{n,i}(y)\) . The mapping \(s_n\) of \(\mathbf{F}\) in \(\mathbf{E}\) is continuous; as the open balls are convex in \(\mathbf{E}\) and in \(\mathbf{F}\) , the mapping \(s_n\) satisfies the inequalities (2) and (3) for all \(y \in \mathbf{F}\) .

From inequality (3) the mappings \(s_n: \mathrm{F} \to \mathrm{E}\) form a Cauchy sequence, for uniform convergence. As \(\mathrm{E}\) is complete, the sequence \((s_n)_{n \in \mathbb{N}}\) converges uniformly to a continuous mapping \(s: \mathrm{F} \to \mathrm{E}\) (GT, X, § 1.6); formula (2) shows that \(u \circ s\) is the identity mapping of \(\mathrm{F}\) , thus \(s\) is a continuous section of \(u\) .

COROLLARY. --- If L is a compact set in F, then there exists a compact set K in E such that \(u(\mathbf{K}) = \mathbf{L}\) .

It is sufficient to put K = s(L), where s is a continuous section of u.

Remarks. --- 1) The corollary to prop. 12 can also be deduced from th. 1 of I, p.~17 and prop. 18 of GT, IX, § 2.10.

\begin{enumerate}
\def\labelenumi{\arabic{enumi})}
\setcounter{enumi}{1}
\tightlist
\item
  We keep the notations of prop. 12. Let \(p\) be a continuous semi-norm on \(E\) ;
\end{enumerate}

for all \(y \in \mathbf{F}\) , put \(q(y) = \inf_{u(x) = y} p(x)\) , so that \(q\) is a continuous semi-norm on \(\mathbf{F}\) (II, p.~4). Let \(\phi\) be a lower semi-continuous mapping of \(F\) in the interval \([0, +\infty[\) of \(\overline{\mathbf{R}}\) . We show that there exists a continuous section \(s\) of \(u\) such that \(p \circ s < q + \phi\) .

Let \(s_0\) be a continuous section of \(u\) (prop. 12) and \(\mathbf{N}\) the kernel of \(u\) . Let \(y_0 \in \mathbf{F}\) , then there exists \(z_0 \in \mathbf{N}\) such that \(p(s_0(y_0) + z_0) < q(y_0) + \phi(y_0)\) . There exists an open neighbourhood \(\mathbf{W}\) of \(y_0\) in \(\mathbf{F}\) such that \(p(s_0(y) + z_0) < q(y) + \phi(y)\) for all \(y \in \mathbf{W}\) . Hence there is an open covering \((\mathbf{W}_i)_{i \in \mathbf{I}}\) of \(\mathbf{F}\) and constant mappings \(t_i: \mathbf{F} \to \mathbf{N}\) such that \(p(s_0(y) + t_i(y)) < q(y) + \phi(y)\) for all \(y \in \mathbf{W}_i\) . As \(\mathbf{F}\) is metrisable, there exists a locally finite continuous partition of unity subordinated to the covering \((\mathbf{W}_i)_{i \in \mathbf{I}}\) , say \((g_i)_{i \in \mathbf{I}}\) (GT, IX, § 4.5, th. 4 and § 4.4, cor. 1). The mapping \(s\) of \(\mathbf{F}\) in \(\mathbf{E}\) defined by \(s(y) = s_0(y) + \sum_{i \in \mathbf{I}} g_i(y). t_i(y)\) fulfils the stated conditions.

